\documentclass[11pt,a4paper]{article}
\usepackage[T1]{fontenc}
\usepackage[utf8]{inputenc}
\usepackage{amsmath,amssymb,amsthm}
\usepackage[margin=1in]{geometry}
\usepackage[colorlinks=true,linkcolor=blue,urlcolor=blue]{hyperref}
\theoremstyle{plain}
\newtheorem{theorem}{Theorem}
\newtheorem{lemma}[theorem]{Lemma}
\newtheorem{proposition}[theorem]{Proposition}
\newtheorem{corollary}[theorem]{Corollary}
\theoremstyle{definition}
\newtheorem{remark}[theorem]{Remark}
% A quoted conjecture can be a single hundred-term formula whose terms are thirty-digit
% numbers. TeX cannot hyphenate those, so without a little stretch every such quote runs off
% the page; the linked-a-file papers are all of that shape.
\setlength{\emergencystretch}{3em}

\title{A further conjecture on OEIS A200787,\\ \large settled by the proof already in hand}
\author{Adrian Perez Fontelles\\ \small Independent researcher}
\date{15 September 2026}
\begin{document}
\maketitle

\begin{abstract}
OEIS A200787 carries more than one conjecture. One of them has a proof; this paper settles a
DIFFERENT one, and it needs no new model and no new engine -- only the linear recurrence already
established for the sequence. A claim about a C-finite sequence whose own characteristic
polynomial is a multiple of the established one is decided by polynomial divisibility, and where
the claim is an identity rather than a recurrence, by divisibility together with a finite window
of agreement.
\end{abstract}

\noindent\small 2020 Mathematics Subject Classification. 11B37, 05A15.\normalsize

\section{The premise}

The entry's first terms are $32, 216, 840, 2425, 5796, 12152, 23136, 40905,\dots$, with offset $1$. The sequence satisfies
\[
  a(n)=6\,a(n-1) - 15\,a(n-2) + 20\,a(n-3) - 15\,a(n-4) + 6\,a(n-5) - 1\,a(n-6) ,
\]
a linear recurrence of order $6$ with characteristic polynomial
\[
  q(x)=x^{6} - 6 x^{5} + 15 x^{4} - 20 x^{3} + 15 x^{2} - 6 x + 1 .
\]
This recurrence is proved in this project; it is the premise of everything below and nothing else is assumed.
It was confirmed against every term the entry publishes from the index where it begins to hold,
with at least $6+2$ further terms after that index.

\section{The further conjecture}

Separately from the conjecture already settled, the entry records

\begin{quote}\raggedright a(n) = (1 + n)\^{}2*(12 + 44*n + 33*n\^{}2 + 7*n\^{}3) / 12 .\end{quote}

The further claim is a CLOSED FORM, and for that the divisibility is not the whole argument. It
gives that $a$ and the closed form $f$ satisfy the same linear recurrence; it does not say they
are the same solution of it. Write $p$ for the annihilator of $f$,
\[
  p(x)=x^{6} - 6 x^{5} + 15 x^{4} - 20 x^{3} + 15 x^{2} - 6 x + 1 ,
\]
of order $D=6$. Since $q\mid p$ and $a$ satisfies the recurrence of $q$, $a$ satisfies that of
$p$ as well; $f$ satisfies it by construction. Two solutions of a linear recurrence of order $D$
with nonzero leading coefficient agree everywhere as soon as they agree at $D$ consecutive
indices, and $a$ and $f$ were checked to agree at $n=7,\dots,12$,
which lies inside the terms the entry publishes and past the index from which $a$ provably
satisfies $p$. Hence $a(n)=f(n)$ for every $n$ in the claimed range.

\begin{theorem}
The conjecture displayed above is true.
\end{theorem}

\section{What is and is not claimed here}

This is a second result on an entry that already has one, and it is a second result only because
it is a DIFFERENT statement. A generating function whose denominator's reciprocal is exactly $q$
says precisely what the proved recurrence says, in another notation; such restatements are
identified and are not counted. The conjecture settled here is not of that kind.

\end{document}
